\section{Conclusion}

This paper presents \sysname{}, a Processing-in-Memory (PIM)
accelerator designed to efficiently serve both transformer and
post-transformer LLMs under the increasing demands of long-context,
high-throughput inference.
%
Through detailed workload characterization, we identify that
\emph{state update}--a central operation in post-transformer
models--shares similar memory bandwidth bottlenecks with attention in
transformer-based models, motivating a unified acceleration approach.
%
\sysname{} addresses these challenges by combining in-memory
computation with quantized execution, co-designing its architecture
around two key principles: (1) maximizing hardware resource sharing
to reduce area cost, and (2) selecting Pareto-optimal quantization
formats for efficient and accurate execution.
%
Our evaluation shows that fine-grained access interleaving and
MX-based state update engines enable \sysname{} to deliver significant
improvements in throughput and area efficiency.
%
These results suggest \sysname{}'s potential as a practical and
generalizable solution for next-generation LLM serving infrastructure.
